\documentclass[../DoAn.tex]{subfiles}
\begin{document}

% ============================================================
%  CHƯƠNG 4 – PHÂN TÍCH THIẾT KẾ, TRIỂN KHAI VÀ ĐÁNH GIÁ
% ============================================================

\section{Thiết kế kiến trúc}

\subsection{Lựa chọn kiến trúc phần mềm}

Hệ thống được xây dựng theo kiến trúc phân tầng kết hợp với kiến trúc thành phần
(layered + component-based), đặc trưng của các dự án trên nền tảng Unity Engine, gồm
năm tầng theo chiều phụ thuộc từ dưới lên. Tầng dữ liệu ở dưới cùng chứa tệp bản đồ định
dạng \texttt{.map} chuẩn MovingAI MAPF benchmark và các \texttt{ScriptableObject} lưu
thông số vật lý xe tăng và đạn, hoàn toàn độc lập với logic gameplay. Trên đó là tầng lõi
MAPF gồm các lớp tìm đường và điều phối đa tác nhân (\texttt{GridAStarPathfinder},
\texttt{GridEnemyAgent}, \texttt{GridEnemyAgentPIBT}, \texttt{GridEnemyAgentPIBT\_TCP}),
không phụ thuộc vào giao diện người dùng. Tầng bootstrap với các lớp \texttt{MapLoader},
\texttt{MapTankTestBootstrap} và \texttt{MapScenarioBootstrap} chịu trách nhiệm khởi tạo
bản đồ, spawn nhân vật và kết nối tầng lõi với tầng vật lý Unity. Tầng gameplay gồm
\texttt{TankController}, \texttt{TankMover} và \texttt{Turret} xử lý vật lý, điều khiển
và bắn đạn cho cả agent AI lẫn người chơi. Trên cùng là tầng backtest với
\texttt{BacktestRunner}, \texttt{BacktestConfigUI} và \texttt{BacktestMode} điều khiển
kịch bản tự động, thu thập chỉ số và xuất CSV.

Việc giữ tầng lõi MAPF tách biệt với giao diện và vật lý giúp thay thế thuật toán
(A*, PIBT C\#, PIBT C++/TCP) mà không thay đổi phần còn lại của hệ thống.

\subsection{Thiết kế tổng quan}

Hình~\ref{fig:arch_packages} mô tả biểu đồ gói UML của hệ thống.
Các gói được tổ chức thành ba nhóm chính theo chiều dọc:
(i)~nhóm hạ tầng bản đồ và tệp dữ liệu nằm ở tầng thấp nhất;
(ii)~nhóm lõi MAPF và điều khiển xe tăng nằm ở tầng giữa;
(iii)~nhóm bootstrap, backtest và giao diện người dùng nằm ở tầng trên cùng.
Chiều phụ thuộc chỉ đi từ trên xuống dưới, không có vòng phụ thuộc ngược.

\begin{figure}[htbp]
    \centering
    \includegraphics[width=0.9\textwidth]{Hinhve/arch_packages.pdf}
    \caption{Biểu đồ gói UML của hệ thống Unity MAPF}
    \label{fig:arch_packages}
\end{figure}

% ============================================================
\section{Thiết kế chi tiết}

\subsection{Thiết kế giao diện}

Ứng dụng chạy trên màn hình độ phân giải tối thiểu 1280\,×\,720, tỷ lệ 16:9, và hỗ trợ
xuất bản WebGL để chạy trên trình duyệt. Giao diện gồm ba màn hình chính. Menu chính cho
phép người dùng lựa chọn chế độ chơi đơn, chơi mạng LAN hoặc chạy backtest. Giao diện
gameplay sử dụng góc nhìn từ trên xuống (top-down 2D), kèm bảng HUD hiển thị điểm máu
Eagle Base, thời gian và số lần tái hoạch định. Giao diện cấu hình backtest cho phép chọn
bản đồ, thuật toán, số lần lặp và bật/tắt chướng ngại động, đồng thời hiển thị thanh tiến
độ và log trực tiếp.

\subsection{Thiết kế các lớp MAPF chính}

Hệ thống AI MAPF được triển khai trong ba biến thể song song nhau, dùng chung tầng điều
khiển xe tăng, như mô tả trong biểu đồ lớp ở Hình~\ref{fig:class_pathfinding}. Biến thể
thứ nhất là \texttt{GridEnemyAgent} kết hợp \texttt{GridAStarPathfinder}, cài đặt thuật
toán A* với heuristic Manhattan trên lưới 4 láng giềng và tái hoạch định định kỳ theo
\texttt{replanInterval}; mỗi agent hoạt động độc lập, không có cơ chế phối hợp. Biến thể
thứ hai là \texttt{GridEnemyAgentPIBT}, cài đặt thuật toán PIBT trực tiếp trong Unity~C\#,
trong đó tại mỗi bước toàn bộ agent chạy đồng thời qua một bộ giải PIBT dùng chung để
đảm bảo không va chạm ở lớp lập kế hoạch. Biến thể thứ ba là \texttt{GridEnemyAgentPIBT\_TCP},
gửi vị trí các agent qua TCP tới máy chủ C++ (PIBT/EPIBT), nhận lại ô mục tiêu tiếp theo
cho từng agent rồi điều khiển vật lý tương tự hai biến thể trên.

\begin{figure}[htbp]
    \centering
    \includegraphics[width=\textwidth]{Hinhve/class_pathfinding.pdf}
    \caption{Biểu đồ lớp các thành phần tìm đường và điều phối agent}
    \label{fig:class_pathfinding}
\end{figure}

Cả ba biến thể đều gọi \texttt{TankController.HandleMoveBody()} và
\texttt{HandleTurretMovement()} -- cùng giao diện mà người chơi sử dụng --
nên hành vi vật lý nhất quán khi so sánh thuật toán.

Hình~\ref{fig:mapf_data_flow} mô tả luồng dữ liệu từ tệp bản đồ đến chuyển động
thực tế của agent trong một chu kỳ tái hoạch định.

\begin{figure}[htbp]
    \centering
    \includegraphics[width=0.88\textwidth]{Hinhve/mapf_data_flow.pdf}
    \caption{Luồng dữ liệu MAPF: từ tệp bản đồ đến chuyển động agent}
    \label{fig:mapf_data_flow}
\end{figure}

\subsection{Thiết kế cơ sở dữ liệu}

Dự án không sử dụng cơ sở dữ liệu quan hệ mà lưu trữ dữ liệu dưới hai dạng tệp. Tệp bản
đồ \texttt{.map} theo định dạng văn bản chuẩn MovingAI MAPF benchmark, với header gồm
\texttt{type}, \texttt{height}, \texttt{width}, \texttt{map} và phần thân là ma trận ký
tự (\texttt{.} cho ô đi được, \texttt{@} cho tường), không thay đổi trong thời gian chạy.
Tệp CSV kết quả backtest được ghi sau mỗi đợt chạy, gồm \texttt{backtest\_summary\_*.csv}
lưu một dòng cho mỗi run và \texttt{backtest\_agents\_*.csv} lưu một dòng cho mỗi agent
trong mỗi run.

\subsection{Thiết kế luồng mạng nhiều người chơi}

Chế độ chơi mạng LAN dựa trên framework Fish-Net theo mô hình một máy chủ và nhiều client.
Quá trình kết nối gồm hai giai đoạn. Giai đoạn tạo và tham gia phòng được mô tả trong
Hình~\ref{fig:seq_lan_create}, trong đó máy chủ khởi tạo phòng và các client lần lượt
tham gia. Giai đoạn sẵn sàng và bắt đầu ván chơi được mô tả trong
Hình~\ref{fig:seq_lan_start}, trong đó hệ thống chờ tất cả người chơi sẵn sàng trước khi
đồng bộ bản đồ và spawn xe tăng cho mọi máy.

\begin{figure}[htbp]
    \centering
    \includegraphics[width=0.95\textwidth]{Hinhve/seq_lan_create_room.pdf}
    \caption{Biểu đồ trình tự tạo và tham gia phòng LAN}
    \label{fig:seq_lan_create}
\end{figure}

\begin{figure}[htbp]
    \centering
    \includegraphics[width=0.95\textwidth]{Hinhve/seq_lan_ready_start.pdf}
    \caption{Biểu đồ trình tự sẵn sàng và bắt đầu ván chơi LAN}
    \label{fig:seq_lan_start}
\end{figure}

% ============================================================
\section{Xây dựng ứng dụng}

\subsection{Công cụ và thư viện sử dụng}

Bảng~\ref{tab:tools} liệt kê các công cụ, thư viện và phiên bản sử dụng trong dự án.

\begin{table}[htbp]
\centering
\caption{Công cụ và thư viện sử dụng}
\label{tab:tools}
\begin{tabular}{lll}
\hline
\textbf{Mục đích}               & \textbf{Công cụ / Thư viện}          & \textbf{Phiên bản} \\ \hline
Engine trò chơi                 & Unity Engine (URP 2D)                & 6000.3.10f1        \\
Ngôn ngữ lập trình              & C\#                                  & .NET 8 (IL2CPP)    \\
Thuật toán MAPF phía server     & PIBT/EPIBT (C++, WSL Ubuntu)         & commit 3f89632     \\
Giao tiếp Unity -- C++ server   & TCP socket (127.0.0.1:7777)          & --                 \\
Chuẩn bản đồ thực nghiệm       & MovingAI MAPF benchmark              & 2019               \\
Môi trường phát triển           & Visual Studio 2022 + Rider           & 17.x / 2024.x      \\
Kiểm soát phiên bản             & Git                                  & 2.44               \\
Xuất / vẽ biểu đồ              & Python 3.14 + matplotlib 3.11        & --                 \\
\hline
\end{tabular}
\end{table}

\subsection{Kết quả đạt được}

Bảng~\ref{tab:codestats} tổng hợp thống kê mã nguồn C\# của hệ thống.

\begin{table}[htbp]
\centering
\caption{Thống kê mã nguồn C\# (tháng 6/2026)}
\label{tab:codestats}
\begin{tabular}{lrr}
\hline
\textbf{Mô-đun}                             & \textbf{Tệp} & \textbf{Dòng lệnh} \\ \hline
Lõi MAPF AI (A*, PIBT, TCP client)          &  8           & 3\,572             \\
Bootstrap và cấu hình kịch bản              &  6           & 3\,617             \\
Backtest (runner, UI, spawner)              &  4           & 1\,820             \\
Gameplay (controller, mover, turret, đạn)  & 12           & 2\,648             \\
Lobby và giao diện mạng LAN                &  9           & 3\,460             \\
Các lớp hỗ trợ (pool, damage, util)        & 38           & 3\,006             \\ \hline
\textbf{Tổng cộng}                          & \textbf{77}  & \textbf{18\,123}   \\ \hline
\end{tabular}
\end{table}

Về mặt phần mềm, đồ án đã hoàn thành một số kết quả cụ thể. Hệ thống đọc và dựng bản đồ
động từ tệp \texttt{.map} chuẩn MAPF, hỗ trợ năm bản đồ kích thước từ 32\,×\,32 đến
251\,×\,180 ô. Ba thuật toán tìm đường đa tác nhân gồm A*, PIBT~C\# và PIBT~C++/TCP hoạt
động ổn định trên cùng năm bản đồ với sáu agent đồng thời. Hệ thống backtest tự động chạy
được 90 kịch bản cho mỗi điều kiện môi trường và xuất kết quả ra CSV cùng biểu đồ. Chế độ
chướng ngại động hoạt động thông qua \texttt{DynamicObstacleSpawner}, ép agent phải tái
hoạch định ngay lập tức khi đường đi bị chặn. Ngoài ra, hệ thống còn hỗ trợ chế độ nhiều
người chơi qua mạng LAN dựa trên framework Fish-Net.

\subsection{Minh họa các chức năng chính}

Phần thực nghiệm và đánh giá (Mục~\ref{sec:experiment}) sẽ minh họa chi tiết hành vi của
các thuật toán qua biểu đồ định lượng. Về mặt trực quan, hệ thống có một số đặc điểm đáng
chú ý. Toàn bộ tường, vật thể và thùng cản đường được xây dựng từ tệp bản đồ tại thời
điểm chạy thay vì bake sẵn vào scene. Khi người dùng chuyển thuật toán giữa A*, PIBT và
PIBT TCP, toàn bộ kịch bản được spawn lại từ đầu mà không cần đóng mở Unity Editor. Trong
chế độ backtest, màn hình hiển thị log tiến độ theo thời gian thực, bao gồm số lần replan,
trạng thái kết nối TCP server và thời gian còn lại của mỗi run.

% ============================================================
\section{Kiểm thử và đánh giá}
\label{sec:experiment}

\subsection{Cấu hình thực nghiệm}

Thực nghiệm được thực hiện theo phương pháp backtest tự động (automated play-out):
agent AI kiểm soát tất cả xe tăng địch, hệ thống ghi lại số chỉ số sau mỗi kịch bản
mà không có sự can thiệp của người dùng.

Bảng~\ref{tab:exp_config} mô tả cấu hình đầy đủ của thực nghiệm.

\begin{table}[htbp]
\centering
\caption{Cấu hình thực nghiệm}
\label{tab:exp_config}
\begin{tabular}{ll}
\hline
\textbf{Tham số}          & \textbf{Giá trị}                                               \\ \hline
Bản đồ                    & Alpha32, Mansion, Chantry, Gallows, Maze128                    \\
Thuật toán so sánh        & A* (baseline), PIBT C\#, PIBT C++/TCP                         \\
Số agent (xe tăng địch)   & 6                                                              \\
Số lần lặp mỗi tổ hợp     & 6                                                              \\
Timeout mỗi run           & 120\,giây (Maze128 chạm trần nội bộ \(\approx\)180\,giây)     \\
Tổng số run / bộ          & 5 bản đồ \(\times\) 3 thuật toán \(\times\) 6 = 90 run        \\
Môi trường tĩnh           & Không có chướng ngại di chuyển                                \\
Môi trường động           & \texttt{DynamicObstacleSpawner}: spawn/despawn thùng          \\
                          & ngay trên đường đi của agent, \texttt{crateLifetime}\,=\,8\,s \\
\hline
\end{tabular}
\end{table}

Hình~\ref{fig:backtest_workflow} mô tả quy trình backtest từ lúc khởi động đến khi
xuất kết quả.

\begin{figure}[htbp]
    \centering
    \includegraphics[width=0.85\textwidth]{Hinhve/backtest_workflow.pdf}
    \caption{Quy trình thực nghiệm backtest tự động}
    \label{fig:backtest_workflow}
\end{figure}

Các chỉ số được thu thập bao gồm: thời gian hoàn thành kịch bản (\textit{Duration}),
tổng số lần tái hoạch định đường đi (\textit{TotalReplans}), số lần phục hồi sau kẹt
(\textit{TotalRecoveries}), tổng số phát bắn (\textit{TotalShots}), tổng số ô đã đi qua
(\textit{TotalCells}), điểm máu Eagle còn lại (\textit{EagleHP}), và kết quả kịch bản
(\textit{Outcome}: \texttt{EagleDestroyed} hoặc \texttt{Timeout}).

\subsection{Kết quả thực nghiệm -- Môi trường tĩnh}

Bảng~\ref{tab:static_results} tổng hợp kết quả trung bình trên 6 lần lặp
trong môi trường tĩnh (không có chướng ngại động).

\begin{table}[htbp]
\centering
\caption{Kết quả thực nghiệm môi trường tĩnh (trung bình 6 lần lặp, 6 agent)}
\label{tab:static_results}
\small
\begin{tabular}{llrrrl}
\hline
\textbf{Bản đồ} & \textbf{Thuật toán} & \textbf{Thời gian (s)} & \textbf{Replan TB} & \textbf{Cells TB} & \textbf{Kết quả} \\ \hline
Alpha32  & A*            & 35,1 &   159  &    64  & EagleDestroyed \\
Alpha32  & PIBT C\#      & 32,3 &   137  &    57  & EagleDestroyed \\
Alpha32  & PIBT C++/TCP  & 35,5 &   137  &   104  & EagleDestroyed \\ \hline
Mansion  & A*            & 98,2 &   672  &   324  & EagleDestroyed \\
Mansion  & PIBT C\#      & 103,3 & 3\,408 &   304  & EagleDestroyed \\
Mansion  & PIBT C++/TCP  & 102,7 & 3\,408 &   472  & EagleDestroyed \\ \hline
Chantry  & A*            & 130,8 &   936  &   458  & EagleDestroyed \\
Chantry  & PIBT C\#      & 127,5 & 3\,514 &   394  & EagleDestroyed \\
Chantry  & PIBT C++/TCP  & 134,4 & 3\,514 &   664  & EagleDestroyed \\ \hline
Gallows  & A*            & 112,4 &   788  &   391  & EagleDestroyed \\
Gallows  & PIBT C\#      & 121,9 & 4\,769 &   388  & EagleDestroyed \\
Gallows  & PIBT C++/TCP  & 115,8 & 4\,769 &   561  & EagleDestroyed \\ \hline
Maze128  & A*            & 180,0 & 1\,439 &   721  & Timeout        \\
Maze128  & PIBT C\#      & 180,0 & 1\,440 &   657  & Timeout        \\
Maze128  & PIBT C++/TCP  & 180,0 & 1\,440 &   990  & Timeout        \\ \hline
\end{tabular}
\end{table}

Hình~\ref{fig:duration_static} và Hình~\ref{fig:replan_static} lần lượt biểu diễn
thời gian hoàn thành trung bình và tổng số lần tái hoạch định trên 5 bản đồ.

\begin{figure}[htbp]
    \centering
    \includegraphics[width=0.92\textwidth]{Hinhve/backtest_duration_static.pdf}
    \caption{Thời gian hoàn thành trung bình theo bản đồ (môi trường tĩnh)}
    \label{fig:duration_static}
\end{figure}

\begin{figure}[htbp]
    \centering
    \includegraphics[width=0.92\textwidth]{Hinhve/backtest_replan_static.pdf}
    \caption{Tổng số lần tái hoạch định trung bình theo bản đồ, thang logarit (môi trường tĩnh)}
    \label{fig:replan_static}
\end{figure}

\subsection{Kết quả thực nghiệm -- Môi trường có chướng ngại động}

Bảng~\ref{tab:dynamic_results} tổng hợp kết quả khi bật chướng ngại động
(thùng xuất hiện ngay trên đường đi của agent trong 8~giây).

\begin{table}[htbp]
\centering
\caption{Kết quả thực nghiệm môi trường có chướng ngại động (trung bình 6 lần lặp, 6 agent)}
\label{tab:dynamic_results}
\small
\begin{tabular}{llrrrl}
\hline
\textbf{Bản đồ} & \textbf{Thuật toán} & \textbf{Thời gian (s)} & \textbf{Replan TB} & \textbf{Cells TB} & \textbf{Kết quả} \\ \hline
Alpha32  & A*            &  40,4 &   238  &    78  & EagleDestroyed \\
Alpha32  & PIBT C\#      &  37,2 &   178  &    69  & EagleDestroyed \\
Alpha32  & PIBT C++/TCP  &  40,8 &   178  &   124  & EagleDestroyed \\ \hline
Mansion  & A*            & 112,9 & 1\,007 &   389  & EagleDestroyed \\
Mansion  & PIBT C\#      & 118,8 & 4\,430 &   365  & EagleDestroyed \\
Mansion  & PIBT C++/TCP  & 118,2 & 4\,430 &   566  & EagleDestroyed \\ \hline
Chantry  & A*            & 150,4 & 1\,403 &   550  & EagleDestroyed \\
Chantry  & PIBT C\#      & 146,6 & 4\,569 &   474  & EagleDestroyed \\
Chantry  & PIBT C++/TCP  & 154,6 & 4\,569 &   798  & EagleDestroyed \\ \hline
Gallows  & A*            & 129,3 & 1\,182 &   469  & EagleDestroyed \\
Gallows  & PIBT C\#      & 140,2 & 6\,201 &   465  & EagleDestroyed \\
Gallows  & PIBT C++/TCP  & 133,2 & 6\,201 &   673  & EagleDestroyed \\ \hline
Maze128  & A*            & 180,0 & 2\,158 &   865  & Timeout        \\
Maze128  & PIBT C\#      & 180,0 & 1\,871 &   788  & Timeout        \\
Maze128  & PIBT C++/TCP  & 180,0 & 1\,871 & 1\,188 & Timeout        \\ \hline
\end{tabular}
\end{table}

Hình~\ref{fig:static_vs_dynamic} so sánh tổng số lần tái hoạch định giữa hai điều kiện môi trường.

\begin{figure}[htbp]
    \centering
    \includegraphics[width=0.88\textwidth]{Hinhve/backtest_static_vs_dynamic.pdf}
    \caption{So sánh số lần tái hoạch định giữa môi trường tĩnh và có chướng ngại động}
    \label{fig:static_vs_dynamic}
\end{figure}

\subsection{Nhận xét và đánh giá}

Từ kết quả trong Bảng~\ref{tab:static_results} và Bảng~\ref{tab:dynamic_results}, có thể
rút ra một số nhận xét. Về thời gian hoàn thành, ba thuật toán cho kết quả xấp xỉ nhau
trên các bản đồ mở như Alpha32, Mansion và Gallows, với chênh lệch lớn nhất không quá
10\%, cho thấy hành vi cuối cùng là phá hủy Eagle Base tương đương nhau giữa các thuật toán.

Về số lần tái hoạch định, hai biến thể PIBT có số replan cao hơn A* khoảng năm đến tám
lần trên các bản đồ hẹp như Mansion, Chantry và Gallows. Nguyên nhân nằm ở chỗ cơ chế
phối hợp của PIBT đếm mỗi bước phân giải xung đột là một lần replan, trong khi A* chỉ
replan khi đường bị chặn hoặc hết chu kỳ định kỳ. Hai giá trị replan của PIBT~C\# và
PIBT~C++/TCP bằng nhau vì cả hai cùng chạy một thuật toán PIBT.

Về số ô đã đi qua, PIBT~C++/TCP đi nhiều ô hơn PIBT~C\# trên mọi bản đồ với chênh lệch từ
20\% đến 70\%. Nguyên nhân là độ trễ TCP giữa Unity và server C++ khiến agent nhận lệnh
chậm hơn một tick, dẫn đến nhiều bước chờ và chuyển hướng thừa.

Về tác động của chướng ngại động, số lần replan tăng rõ rệt khi bật chướng ngại, trung
bình khoảng 50\% với A* và 30\% với PIBT, trong khi thời gian hoàn thành tăng khoảng 15\%
do agent phải tính đường vòng. Kết quả này xác nhận cả ba thuật toán đều có khả năng tái
hoạch định khi môi trường thay đổi. Riêng bản đồ mê cung Maze128 kích thước 128$\times$128
vượt quá ngưỡng timeout với cả ba thuật toán ở cả hai điều kiện môi trường; đây là kết
quả hợp lệ, phản ánh giới hạn về tốc độ tiếp cận mục tiêu trên bản đồ mê cung lớn với
sáu agent.

% ============================================================
\section{Triển khai}

Hệ thống được triển khai theo hai hình thức. Hình thức thứ nhất là chạy trực tiếp trong
Unity Editor phục vụ phát triển và thực nghiệm: người dùng mở scene
\texttt{Assets/Scenes/MapF\_TankTest.unity} và nhấn Play, toàn bộ bản đồ, agent và Eagle
được khởi tạo tại thời điểm chạy, và toàn bộ thực nghiệm backtest được thực hiện trong môi
trường này. Hình thức thứ hai là triển khai trực tuyến qua WebGL: ứng dụng được xuất sang
WebGL và phục vụ qua GCP Nginx/HTTPS, cho phép truy cập từ trình duyệt mà không cần cài
đặt thêm phần mềm; tuy nhiên chế độ backtest không khả dụng trên WebGL do giới hạn truy
xuất tệp của trình duyệt.

Máy chủ PIBT C++ (\texttt{Server-PIBT-TeamNoMan-sSky}) chạy độc lập trên WSL/Ubuntu
và lắng nghe kết nối TCP tại cổng 7777. Khi chạy thực nghiệm, Unity kết nối tới
server này qua localhost; kết quả backtest không bao gồm độ trễ mạng thực tế.

\end{document}
